% !TEX root = ../notes_sovereignai.tex
% Chapter 2: The Engine. First boot, OS, driver, Docker, Ollama, Open WebUI.
%
% This chapter owns everything from the first power-on onward. Firmware setup and
% the enumeration probe live here, not in the hardware chapter: they are the first
% thing software does to the rig, and they are what separates a hardware fault from
% a software one before there is any software to blame.
%%%%%%%%%%%%%%%%%%%%%%%%%%%%%%%%%%%%%%%%%%%%%%%%%%%%%%%%%%%%%%%%%%%%%%

\index{Ollama}\index{Open WebUI}\index{Docker}\index{CUDA}\index{Ubuntu Server}\index{NVIDIA driver}
\chapter{The Engine}
\label{ch:engine}
\printchapterversion{02_engine}
\newthought{Silicon, Meet Software.}
\vspace{1.5em}

% ==============================================================================================
% THE WHY
% ==============================================================================================

\section{The Why}\index{inference}\index{container runtime}

\newthought{The rig is assembled}, wired, and has never been switched on. A graphics card is
only potential until software puts a model on it and hands you a way to ask
questions. This chapter powers the machine for the first time, proves every part
is really there, and turns the metal into a served model the whole family can talk
to, with the least machinery that works: firmware, an operating system, a driver, a
container runtime, and one program that loads models onto the GPU.

\newthought{That program is Ollama.}\index{Ollama!install} It is the shortest path to local
inference: one command installs it, it finds the GPU on its own, and it serves an
API\index{OpenAI-compatible API} the rest of the world already speaks. In front of it we put
Open WebUI, a private chat in the browser, so the family has something to use on
day one and you have an account system to build on later.

\marginnote{\newthought{Why not vLLM yet.}\index{vLLM} vLLM wins when many people hit the
box at once, and it is also what answers a narrow slot: cards two to four sit on a
single PCIe lane, where tensor parallelism\index{tensor parallelism} would push activations across the
link on every token. vLLM need not split a model that way. A replica per card, or a
split by layer with pipeline parallelism\index{pipeline parallelism}, sends almost nothing over the link,
so lane width stops mattering and the answer stays in software rather than in a
server board. On one card serving a household, Ollama's one command beats vLLM's
tuning, and both speak the same OpenAI-compatible API, so the swap is later and cheap.}

\newthought{By the end of this chapter} you will have:
\begin{itemize}
    \item set five firmware values and proved the hardware enumerates before an
    operating system touched the drives;
    \item put Ubuntu Server, the NVIDIA driver, and GPU-aware Docker on the rig;
    \item run a model on the 3090 with Ollama and confirmed it answers on the API;
    \item stood up Open WebUI and chatted with the model from a browser on your LAN.
\end{itemize}

% ==============================================================================================
% FIRST BOOT
% ==============================================================================================

\section{First Boot}\index{firmware}

\newthought{Power on.} Connect the monitor to the board's own display output, not to
the graphics card, switch the supply on, and press the power button. The fans should
spin up and hold; a half turn and a stop means the supply or the 24~pin is not
seated. Enter setup when the splash appears.

\marginnote{\newthought{Two of these are for later.} Above 4G Decoding and
Resizable BAR only bite once a second card arrives, so set them now, while the
rail is still empty and a later card cannot fail in a way that looks like dead
hardware.}

\newthought{Set five values}\index{Above 4G Decoding}\index{Resizable BAR}\index{IOMMU}\index{integrated graphics}\index{EXPO profile}\index{XMP profile}, then save and exit:

\begin{verbatim}
Above 4G Decoding ............ Enabled
Re-Size BAR Support .......... Enabled
IOMMU ........................ Enabled
Initial Display Output ....... IGD
EXPO / XMP profile ........... Profile 1
\end{verbatim}

\marginnote{\newthought{What they do.} Above 4G Decoding gives the cards address
space beyond the 32~bit boundary, without which a second card fails to map at all.\index{VRAM}
Resizable BAR lets the processor address a card's whole VRAM in one window, which
speeds weight loading. IOMMU keeps device passthrough\index{device passthrough} available. Initial Display
Output on the integrated graphics leaves all 24~GB of every card free for models.}

\newthought{The last of the five is the memory.} A DDR5 kit does not run at the
speed on its box until you say so: leave the EXPO profile off and the modules clock
themselves to the conservative default published by JEDEC\index{JEDEC}, the memory standards
body, rather than the rated 6000 they were bought for. This is the one firmware
value that changes nothing about whether the machine works and everything about how
fast it feeds the cards.

\marginnote{\newthought{Menu names move.} These five sit under different headings
in different firmware revisions of the same board. Use the setup search rather than
trusting a path from a screenshot, and set all five in one visit: a value missed
here surfaces later as a hardware fault that is not one.}

\newthought{Probe before you install anything.}\index{lspci}\index{enumeration} Boot an Ubuntu Server image from a
USB stick and choose to try it without installing. Nothing here writes to the
drives:

\begin{verbatim}
$ lscpu | grep 'Model name'   # AMD Ryzen 5 9600X
$ free -h                     # total ~ 62Gi
$ lspci | grep -i nvidia      # NVIDIA GA102 [RTX 3090]
$ lspci -vv | grep LnkSta:    # Width x16
\end{verbatim}

\marginnote{\newthought{Why this proves anything.} Enumeration is the firmware
naming each device on the bus, and it happens with no driver loaded. A card that
enumerates is seated, powered, and addressable, which separates a hardware fault
from a software one before there is any software to blame.}

\newthought{Four clean answers} clear the way for software: the right processor, the
full memory, an NVIDIA device on the bus, and that device at x16. If the card is
missing, reseat the riser and confirm both eight pin cables. If the width reads x1,
the card is on the wrong slot or the wrong riser. If memory reads half, reseat the
second module. Past this point every failure is a software failure.

\marginnote{\newthought{Link speed, read honestly.}\index{PCIe!link width} \texttt{LnkSta} is the
negotiated state, not the capability, and a card at idle may report a lower speed
until it is put to work. Width is the number to trust here.}

% ==============================================================================================
% LOCAL SERVING
% ==============================================================================================

\section{Local Serving}\index{quantization}\index{model weights}

\newthought{A served model is a stack} four layers deep. The operating system
owns the hardware; the NVIDIA driver\index{NVIDIA driver} lets programs reach the card; the NVIDIA
Container Toolkit\index{NVIDIA Container Toolkit} lets Docker containers reach it too, so every service in this
book ships as a container without fighting over the GPU; and Ollama sits on top,
loading model weights\index{model weights} into the 24~GB of video memory\index{VRAM} and answering requests.

\marginnote{\newthought{Why containers.} Each later service, the chat, the voice
pipeline, image generation, is a container that asks Docker for the GPU\index{Docker!GPU passthrough} and gets
it. The toolkit passes the host driver through, so you install the driver once on
the metal and never again inside a container.}

\newthought{Models come quantized.} A model's weights are shrunk from sixteen bit
floats to four or five bits with little loss\index{quantization}, which is what lets a capable model
fit in 24~GB. Ollama pulls these as ready-to-run files and picks a sensible
quantization for you; a fourteen billion parameter model leaves headroom for
context, a thirty billion parameter mixture-of-experts\index{mixture of experts} model fills the card and is
the ceiling for a single 3090.

\marginnote{\newthought{Pick a tool-caller.}\index{tool calling} Choose from the Qwen3\index{Qwen3} family. Beyond
chat quality, prefer a model that supports tool calling, since Chapter~4's voice
assistant cannot flip a single switch without it. Filter the model library by
\texttt{tools} and ignore the rest.}

\newthought{The API is the point.}\index{Ollama!API} Ollama serves an OpenAI-compatible endpoint at
\texttt{:11434/v1}, the same shape every client in this book speaks. Two cautions
that Chapter~3 acts on: it binds to localhost and it has no authentication, so on
its own it is reachable only from the machine itself, and opening it to the
network without a gateway in front would leave it wide open\cite{ollamaDocs}.

\section{Run a Model}\index{nvidia-smi}

\marginnote{\newthought{Headless is the habit.}\index{SSH} Do all of this over SSH from your
laptop; the rig needs no monitor of its own. Run the GPU steps first, because every
later one assumes the card is visible to software.}

\newthought{Lay the foundation.}\index{Ubuntu Server} Install Ubuntu Server~24.04 from the same USB
stick you probed with, then the driver and a GPU-aware Docker:

\begin{verbatim}
$ sudo ubuntu-drivers install          # NVIDIA driver
$ sudo reboot
$ nvidia-smi                           # the 3090 appears, with its 24 GiB

$ curl -fsSL https://get.docker.com | sh
$ sudo nvidia-ctk runtime configure --runtime=docker
$ sudo systemctl restart docker
$ docker run --rm --gpus all ubuntu nvidia-smi   # the card, now inside a container
\end{verbatim}

\marginnote{\newthought{Split the drive here, not at the shop.}\index{partitioning}\index{Ollama!model directory} Manual partitioning
in the installer gives the system its own partition and leaves the rest as a
separate filesystem for weights, and the board keeps two free M.2 sockets if a
second drive comes later. Either way \texttt{OLLAMA\_MODELS} points at the large
volume, so reinstalling the operating system never touches a model library that
took days to download.}

\newthought{Install the engine and pull a model.}\index{Ollama!model pull}\index{model library} One command installs Ollama,
which detects the card; one more pulls a model and opens a chat:

\begin{verbatim}
$ curl -fsSL https://ollama.com/install.sh | sh
$ ollama pull qwen3:14b
$ ollama run qwen3:14b      # type a question; watch nvidia-smi load the card
\end{verbatim}

\newthought{Confirm it answers on the API}, the interface every later chapter
uses:

\begin{verbatim}
$ curl http://localhost:11434/v1/chat/completions -d '{
    "model": "qwen3:14b",
    "messages": [{"role": "user", "content": "say hello"}]
  }'
\end{verbatim}

\newthought{Put a face on it.}\index{Docker!volume} Run Open WebUI as a container pointed at Ollama,
with a named volume so accounts and history survive restarts:

\begin{verbatim}
$ docker run -d --name open-webui -p 3000:8080 \
    --add-host=host.docker.internal:host-gateway \
    -e OLLAMA_BASE_URL=http://host.docker.internal:11434 \
    -v open-webui:/app/backend/data \
    ghcr.io/open-webui/open-webui:main
\end{verbatim}

\marginnote{\newthought{The first account is admin.}\index{Open WebUI!accounts} Open WebUI's first sign-up
becomes the administrator; create it, then add an account for each member of the
household under the admin panel. That account system is the gateway Chapter~3 hands
keys from.}

\newthought{What you now own} is a served model. From any browser on your home
network, open \texttt{http://rig-address:3000}, sign in, and chat. A family member
two rooms away is talking to a model running on your 3090, with no token metered
and no prompt leaving the house.

% ==============================================================================================
% BRIDGE
% ==============================================================================================

\newthought{The engine runs, but only at home.} It answers on the local network and
nowhere else, and it trusts anyone who can reach it. The next chapter carries it to
your phone, your laptop, and your terminal wherever they are, without opening the
box to everyone.

\marginnote{\newthought{feedback}}
